\chapter{Opis metod i technologii wykorzystanych w pracy}

\section{Język programowania}

W niniejszej pracy wykorzystany został język programowania Python w wersji 3.14.6 \cite{python} jako główne narzędzie do implementacji opracowanego rozwiązania. Wybór ten wynikał z faktu, że jest to stabilna i dobrze przetestowana dystrybucja tego oprogramowania. Pozwala ona na korzystanie z wielu zewnętrznych bibliotek i paczek bez konfliktów związanych z niezgodnością wersji. Oprócz tego z każdą kolejną aktualizacją języka dodawane są do niego kolejne usprawnienia dlatego w miarę możliwości powinno używać się go w jak najnowszej wersji.

\section{Zewnętrzne paczki pythonowe}

W projekcie wykorzystano zestaw zewnętrznych bibliotek języka Python. Ich konkretne wersje zostały określone w pliku \textbf{requirements.txt}, dzięki czemu możliwe jest odtworzenie środowiska wykorzystanego w pracy.

\begin{itemize}

\item \textbf{qiskit 2.5.0} \cite{qiskit}
\item \textbf{qiskit\_algorithms 0.4.0} \cite{qiskit_algorithms}
\item \textbf{networkx 3.6.1} \cite{networkx}
\item \textbf{numpy 2.5.1} \cite{numpy}
\item \textbf{pandas 3.0.3} \cite{pandas}
\item \textbf{matplotlib 3.11.1} \cite{matplotlib}

\end{itemize}
